\documentclass[12pt,a4paper]{article}
\usepackage[T1]{fontenc}
\usepackage{amsmath,amssymb,mathrsfs,tikz,times,pifont}
\usepackage{enumitem}
\usepackage{multicol}
\usepackage{lmodern}
\usetikzlibrary{trees}
\newcommand\circitem[1]{%
\tikz[baseline=(char.base)]{
\node[circle,draw=gray, fill=red!55,
minimum size=1.2em,inner sep=0] (char) {#1};}}
\newcommand\boxitem[1]{%
\tikz[baseline=(char.base)]{
\node[fill=cyan,
minimum size=1.2em,inner sep=0] (char) {#1};}}
\setlist[enumerate,1]{label=\protect\circitem{\arabic*}}
\setlist[enumerate,2]{label=\protect\boxitem{\alph*}}
%%%::::::by chnini ameur :::::::%%%
\everymath{\displaystyle}
\usepackage[left=1cm,right=1cm,top=1cm,bottom=1.7cm]{geometry}
\usepackage[colorlinks=true, linkcolor=blue, urlcolor=blue, citecolor=blue]{hyperref}
\usepackage{array,multirow}
\usepackage[most]{tcolorbox}
\usepackage{varwidth}
\usepackage{float} %pour utiliser l'option [H] qui force l'image à apparaître exactement à l'endroit où elle est placée dans le code.
\tcbuselibrary{skins,hooks}
\usetikzlibrary{patterns}
%%%::::::by chnini ameur :::::::%%%
\newtcolorbox{exa}[2][]{enhanced,breakable,before skip=2mm,after skip=5mm,
colback=yellow!20!white,colframe=black!20!blue,boxrule=0.5mm,
attach boxed title to top left ={xshift=0.6cm,yshift*=1mm-\tcboxedtitleheight},
fonttitle=\bfseries,
title={#2},#1,
% varwidth boxed title*=-3cm,
boxed title style={frame code={
\path[fill=tcbcolback!30!black]
([yshift=-1mm,xshift=-1mm]frame.north west)
arc[start angle=0,end angle=180,radius=1mm]
([yshift=-1mm,xshift=1mm]frame.north east)
arc[start angle=180,end angle=0,radius=1mm];
\path[left color=tcbcolback!60!black,right color = tcbcolback!60!black,
middle color = tcbcolback!80!black]
([xshift=-2mm]frame.north west) -- ([xshift=2mm]frame.north east)
[rounded corners=1mm]-- ([xshift=1mm,yshift=-1mm]frame.north east)
-- (frame.south east) -- (frame.south west)
-- ([xshift=-1mm,yshift=-1mm]frame.north west)
[sharp corners]-- cycle;
},interior engine=empty,
},interior style={top color=yellow!5}}
%%%%%%%%%%%%%%%%%%%%%%%
\usepackage{fancyhdr}
\usepackage{eso-pic}         % Pour ajouter des éléments en arrière-plan
% Commande pour ajouter du texte en arrière-plan
\usepackage{tkz-tab}
\AddToShipoutPicture{
    \AtTextCenter{%
        \makebox[0pt]{\rotatebox{80}{\textcolor[gray]{0.7}{\fontsize{5cm}{5cm}\selectfont PGB}}}
    }
}
\usepackage{lastpage}
\fancyhf{}
\pagestyle{fancy}
\renewcommand{\footrulewidth}{1pt}
\renewcommand{\headrulewidth}{0pt}
\renewcommand{\footruleskip}{10pt}
\fancyfoot[R]{
\color{blue}\ding{45}\ \textbf{2025}
}
\fancyfoot[L]{
\color{blue}\ding{45}\ \textbf{Prof:M. BA}
}
\cfoot{\bf
\thepage /
\pageref{LastPage}}
% Création du compteur pour les exercices
\newcounter{exercice}
\renewcommand{\theexercice}{\arabic{exercice}}  % Définit l'affichage du compteur en chiffres arabes

% Définir la commande \exo
\newcommand{\exo}{\refstepcounter{exercice}\textbf{Exercice \theexercice} }
\begin{document}
\renewcommand{\arraystretch}{1.5}
\renewcommand{\arrayrulewidth}{1.2pt}
\begin{tikzpicture}[overlay,remember picture]
    \node[draw=blue,line width=1.2pt,fill=purple,text=blue,inner sep=3mm,rounded corners,pattern=dots]at ([yshift=-2.5cm]current page.north) {\begingroup\setlength{\fboxsep}{0pt}\colorbox{white}{\begin{tabular}{|*1{>{\centering \arraybackslash}p{0.28\textwidth}} |*2{>{\centering \arraybackslash}p{0.2\textwidth}|} *1{>{\centering \arraybackslash}p{0.19\textwidth}|} }
                \hline
                \multicolumn{3}{|c|}{$\diamond$$\diamond$$\diamond$\ \textbf{Lycée de Dindéfélo}\ $\diamond$$\diamond$$\diamond$ } & \textbf{A.S. : 2025/2026}                                              \\ \hline
                \textbf{Matière: Mathématiques}                                                                                    & \textbf{Niveau : T}\textbf{S2} & \textbf{Date: 05/12/2025} & \textbf{} \\ \hline
                \multicolumn{4}{|c|}{\parbox[c]{10cm}{\begin{center}
                                                                  \textbf{{\Large\sffamily Td Complexe}}
                                                              \end{center}}}                                                                                                        \\ \hline
            \end{tabular}}\endgroup};
\end{tikzpicture}
\vspace{3cm}


\fbox{\textbf{\exo}} [Module et argument]

\fbox{\textbf{\exo}} : [Forme trigonométrique]

\fbox{\textbf{\exo}} [Formes trigonométrique et exponentielle]

\fbox{\textbf{\exo}} [Forme trigonométrique]

\fbox{\textbf{\exo}} [Classique !]

\fbox{\textbf{\exo}} [Racines carrées et Equations]

\fbox{\textbf{\exo}} [Racines $n$-ièmes]
\fbox{\textbf{\exo}} [Racines $n$-ièmes]

On donne $z_0 = 1 - i\sqrt{3}$

\begin{enumerate}
    \item Donner une écriture trigonométrique de $z_0$.
    \item Montrer que $z_0^4 = -8 + 8i\sqrt{3}$.
    \item Résoudre dans $\mathbb{C}$ l'équation $z^4 = 1$.
    \item En déduire les solutions de $(E) : z^4 = -8 + 8i\sqrt{3}$ sous forme algébrique et sous forme trigonométrique.
\end{enumerate}
\fbox{\textbf{\exo}} [Racines $n$-ièmes]

\fbox{\textbf{\exo}}[Linéarisation]

\fbox{\textbf{\exo}} [Lieux géométriques]

\fbox{\textbf{\exo}} [Lieux géométriques]

\fbox{\textbf{\exo}} [Deux méthodes pour un problème]

\fbox{\textbf{\exo}} [Complexe et géométrie]

\fbox{\textbf{\exo}}

\fbox{\textbf{\exo}}

\fbox{\textbf{\exo}}

\fbox{\textbf{\exo}}

\fbox{\textbf{\exo}}

\fbox{\textbf{\exo}}

\fbox{\textbf{\exo}}

\fbox{\textbf{\exo}}

\fbox{\textbf{\exo}}

\fbox{\textbf{\exo}}

\fbox{\textbf{\exo}}

\fbox{\textbf{\exo}}
\section*{\underline{\textbf{Exercice :}} Nombres complexes}

On donne 
\[
a = \sqrt{2 - \sqrt{3}} - i\sqrt{2 + \sqrt{3}}.
\]

\begin{enumerate}
    \item \textbf{Calcul de $a^2$ et du module de $a$}

    Posons 
    \[
    x = \sqrt{2 - \sqrt{3}}, \qquad y = \sqrt{2 + \sqrt{3}}.
    \]
    Alors $a = x - iy$.

    On calcule :
    \[
    a^2 = (x - iy)^2 = x^2 - 2ixy - y^2.
    \]

    Or
    \[
    x^2 = 2 - \sqrt{3}, \qquad y^2 = 2 + \sqrt{3}.
    \]

    Donc
    \[
    a^2 = (2 - \sqrt{3}) - (2 + \sqrt{3}) - 2i\sqrt{(2-\sqrt3)(2+\sqrt3)}
    = -2\sqrt{3} - 2i.
    \]

    Ainsi,
    \[
    \boxed{a^2 = -2\sqrt{3} - 2i.}
    \]

    Le module de $a$ vérifie :
    \[
    |a|^2 = |a^2| = \sqrt{(-2\sqrt3)^2 + (-2)^2} = \sqrt{12 + 4} = 4,
    \]
    d'où
    \[
    \boxed{|a| = 2.}
    \]

    \item \textbf{Forme trigonométrique de $a^2$ et argument de $a$}

    On a
    \[
    a^2 = -2\sqrt{3} - 2i.
    \]

    Son module est $|a^2| = 4$. On écrit
    \[
    a^2 = 4(\cos\theta + i\sin\theta),
    \]
    avec
    \[
    \cos\theta = \frac{-2\sqrt3}{4} = -\frac{\sqrt3}{2}, 
    \qquad
    \sin\theta = \frac{-2}{4} = -\frac{1}{2}.
    \]

    On reconnaît
    \[
    \theta = \frac{7\pi}{6} \quad (\text{ou } -\tfrac{5\pi}{6}).
    \]

    Donc
    \[
    \boxed{a^2 = 4\left(\cos\frac{7\pi}{6} + i\sin\frac{7\pi}{6}\right).}
    \]

    Comme $a^2 = |a|^2 e^{i2\arg(a)}$, on a :
    \[
    2\arg(a) = \frac{7\pi}{6} + 2k\pi
    \quad\Rightarrow\quad
    \arg(a) = \frac{7\pi}{12} + k\pi.
    \]

    Une valeur positive possible est donc :
    \[
    \boxed{\arg(a) = \frac{19\pi}{12}.}
    \]

    \item \textbf{Valeurs exactes des cosinus et sinus}

    On a 
    \[
    a = 2\left(\cos\frac{19\pi}{12} + i\sin\frac{19\pi}{12}\right).
    \]

    Donc
    \[
    \cos\frac{19\pi}{12} = \frac{\Re(a)}{2} = \frac{\sqrt{2-\sqrt3}}{2}, 
    \]
    \[
    \sin\frac{19\pi}{12} = \frac{\Im(a)}{2} = -\frac{\sqrt{2+\sqrt3}}{2}.
    \]

    Or 
    \[
    \frac{19\pi}{12} = \frac{7\pi}{12} + \pi,
    \]
    donc
    \[
    \cos\frac{19\pi}{12} = -\cos\frac{7\pi}{12}, 
    \qquad
    \sin\frac{19\pi}{12} = -\sin\frac{7\pi}{12}.
    \]

    Ainsi :
    \[
    \boxed{
    \begin{aligned}
    \cos\frac{7\pi}{12} &= \frac{\sqrt{2-\sqrt3}}{2},\\
    \sin\frac{7\pi}{12} &= \frac{\sqrt{2+\sqrt3}}{2}.
    \end{aligned}}
    \]

    Enfin,
    \[
    \boxed{
    \begin{aligned}
    \cos\frac{\pi}{12} &= \frac{\sqrt{2+\sqrt3}}{2},\\
    \sin\frac{\pi}{12} &= \frac{\sqrt{2-\sqrt3}}{2}.
    \end{aligned}}
    \]

\end{enumerate}

\fbox{\textbf{\exo}}

\fbox{\textbf{\exo}}(POLYNÔMES COMPLEXES)

\fbox{\textbf{\exo}} (POLYNÔMES COMPLEXES)

\fbox{\textbf{\exo}} (POLYNÔMES COMPLEXES)

\fbox{\textbf{\exo}}

\end{document}
